% sections/11_road_features.tex — V7

\section{V7: Severity in Presence of Road Features}

Four small multiples compare the severity split (High/Low) for accidents at
locations \emph{with} vs.\ \emph{without} Traffic Signals, Crossings,
Junctions, and Stop Signs.

% \figchart[0.95]{v7_road_features}{Road feature presence vs severity split}{fig:v7}
\begin{figure}[H]
  \centering
  \placeholder[6cm]{%
    \textbf{V7 --- Small Multiples: Road Feature vs Severity Split}\\[0.5em]
    4 panels: Traffic Signal / Crossing / Junction / Stop Sign\\[0.4em]
    Save as \texttt{figures/v7\_road\_features.png}%
  }
  \caption{High-severity rate in presence vs.\ absence of road features}
  \label{fig:v7}
\end{figure}

\begin{table}[H]
\centering\small
\begin{tabular}{@{}lccl@{}}
  \toprule
  \hdrrow
  \textbf{Road Feature}
    & \textbf{High Sev.\ (Present)}
    & \textbf{High Sev.\ (Absent)}
    & \textbf{Effect} \\ \midrule
  \posrow Traffic Signal & 9.51\%  & 21.19\% & $+11.68$\% (increases) \\
  \posrow Crossing       & 7.04\%  & 21.04\% & $+14.00$\% (increases) \\
  \negrow Junction       & 26.79\% & 18.88\% & $-7.91$\% (reduces)\\
  \posrow Stop Sign      & 6.22\%  & 19.84\% & $+13.62$\% (increases) \\
  \bottomrule
\end{tabular}
\caption{Road feature presence vs.\ high-severity rate}
\end{table}

\insightbox{Infrastructure with mandatory stopping behaviour (signals, crossings,
stop signs) substantially lowers high-severity rates, while junctions --- which
require yield decisions without mandatory stops --- increase severity.
Improving junction control (adding signals or roundabouts) at high-accident
junctions could meaningfully reduce severity by improving traffic organisation
and reducing road clearance time.}
